\chapter{איומים וחולשות}

\section{סוגי איומים על מודלים של שפה}

כיום, כאשר מודלים של שפה משמשים יותר ויותר בתחומים רגישים כגון בריאות, כספים, וחינוך, סוגיות של בטיחות והיכחשות קריטיים. יש מספר סוגים של איומים שעלול להיתקל בהם בעת השימוש במודלים כאלה:

\begin{enumerate}
    \item \textenglish{Hallucinations}: המודל יוצר מידע שאינו נכון או שלא קיים בנתונים ההדרכה.
    \item \textenglish{Prompt injection}: משתמש זדוני יכול לתחברק את הקלט כדי להכניס הוראות מסתתרות הגורמות למודל להתנהג בצורה לא רצויה.
    \item \textenglish{Data poisoning}: כאשר מישהו מיתח את נתוני ההדרכה כדי להשפיע על התנהגות המודל.
    \item \textenglish{Adversarial examples}: קלט מעוטר בדרך ערמומית כדי לגרום למודל לטעות בדרך מסוימת.
    \item \textenglish{Privacy leakage}: המודל יכול לגלות מידע רגיש שהיה בנתוני ההדרכה.
\end{enumerate}

\section{איומים על פרוטוקול בקרת המודל}

כאשר משתמשים ב-\textenglish{MCP} (פרוטוקול בקרת המודל), נוצרים איומים נוספים שקשורים לתקשורת בין הישום ובין המודל. למשל, אם המודל יקבל הוראה מיישום חיצוני שלא עבר בדיקה, הוא יכול להתנהג בדרך לא בטוחה.

בנוסף, ה-\textenglish{context window} של מודל (גודל הקלט המרבי שהוא יכול לקבל) יכול להיות נוצל בצורה שלילית. מודל כגון \textenglish{Claude} יכול לקבל עד כ-200 אלף טוקנים בקלט, מה שפותח אפשרויות להשפעות בקשר ישירות.

סוגיה משמעותית נוספת היא זו של \textenglish{authorization}: כיצד נוודא שגם כשמודל משתמש בכלים או מידע חיצוני דרך \textenglish{MCP}, הוא עדיין משתמש בהן בדרך המתאימה לרמת הזכויות של משתמש מסוים?

\section{אתגרים של קנה מידה}

ככל שמודלים הופכים יותר גדולים וכוחניים, אתגרים של בקרה הופכים יותר גדולים. מודל כגון \textenglish{GPT-3} עם 175 מיליארד פרמטרים יכול להפוך כל הנוף של סוגיות. האתגרים הוא לא רק לבנות בדיקות טובות, אלא גם לבנות אלה שנשמרים עם הגדלת המודל עצמו.

כמו כן, יש צורך לפתח גישות למידה שלא דורשות מליליוני שעות של בדיקה אנושית, אלא מאמצות שמודלים עצמם יכולים להשתמש בהם לבדיקה ולשיפור שלהם.
